\documentclass{article}
\usepackage{lmodern}
\usepackage{amsmath}
\usepackage{listings}
\usepackage{fullpage}
\usepackage{parskip}
\usepackage{xcolor}
\lstset{
	basicstyle=\ttfamily,
	columns=fullflexible,
	frame=single,
	breaklines=true,
	postbreak=\mbox{\textcolor{red}{$\hookrightarrow$}\space},
}
\begin{document}
\title{Experiment `p\_4\_6\_find\_minimum\_seq' Results}
\author{\today}
\date{}
\maketitle
\textbf{Experiment outcome: }SUCCESS
\\\textbf{Bad responses: }0
\\\textbf{Responses containing}~\texttt{assume}~\textbf{: }0
\\\textbf{Responses containing}~\texttt{expect}~\textbf{: }0
\\\textbf{Resolution attempts: }1
\\\textbf{Hard fails (resolution): }0
\\\textbf{Soft fails (resolution): }0
\\\textbf{Verification attempts: }1
\section*{Problem Specification}
\textbf{Problem name: }p\_4\_6\_find\_minimum\_seq
\\\textbf{Natural language statement: }null
\\\textbf{Method signature: }p\_4\_6\_find\_minimum\_seq(inputs: seq<int>) returns (min: int)
\subsection*{Ensures}
\begin{itemize}
\item \texttt{forall i :: 0 <= i < |inputs| ==> min <= inputs[i]}
\item \texttt{exists i :: 0 <= i < |inputs| \&\& min == inputs[i]}
\end{itemize}
\subsection*{Requires}
\begin{itemize}
\item \texttt{0 < |inputs|}
\end{itemize}
\clearpage
\section*{GenAI interactions}
Below you will find all interactions between the `user' (program) and the `assistant' (OpenAI).
\subsection*{Program $\rightarrow$ GenAI}
\begin{lstlisting}
You are given the following task to perform in Dafny:

Write a method to satisfy the below requirements.

The signature should be:

p_4_6_find_minimum_seq(inputs: seq<int>) returns (min: int)

The method should respect the following contract:

requires 0 < |inputs|, ensures forall i :: 0 <= i < |inputs| ==> min <= inputs[i], ensures exists i :: 0 <= i < |inputs| && min == inputs[i]

Produce and show only the Dafny body of this method, including the curly braces that surround it. Do not show the signature nor contract. You must not use 'assume' nor 'expect' anywhere in your code.
\end{lstlisting}
\subsection*{GenAI $\rightarrow$ Program}
\textbf{System fingerprint: }null
\\\textbf{ID: }chatcmpl-DnMDuoa62x4jvmrDq3n8GCona37NY
\\\textbf{Model: }gpt-5.5-2026-04-23
\\\textbf{Created at: }1780655870
\\\textbf{Finish reason: }stop
\\\textbf{Usage: }322 tokens in, and 1023 tokens out
\begin{lstlisting}
{
  var idx := 1;
  var minIndex := 0;
  min := inputs[0];

  while idx < |inputs|
    invariant 1 <= idx <= |inputs|
    invariant 0 <= minIndex < idx
    invariant min == inputs[minIndex]
    invariant forall i :: 0 <= i < idx ==> min <= inputs[i]
  {
    if inputs[idx] < min {
      min := inputs[idx];
      minIndex := idx;
    }
    idx := idx + 1;
  }
}
\end{lstlisting}
\section*{Final Program}
\begin{lstlisting}
method p_4_6_find_minimum_seq(inputs: seq<int>) returns (min: int)
	requires 0 < |inputs|
	ensures forall i :: 0 <= i < |inputs| ==> min <= inputs[i]
	ensures exists i :: 0 <= i < |inputs| && min == inputs[i]
{
  var idx := 1;
  var minIndex := 0;
  min := inputs[0];

  while idx < |inputs|
    invariant 1 <= idx <= |inputs|
    invariant 0 <= minIndex < idx
    invariant min == inputs[minIndex]
    invariant forall i :: 0 <= i < idx ==> min <= inputs[i]
  {
    if inputs[idx] < min {
      min := inputs[idx];
      minIndex := idx;
    }
    idx := idx + 1;
  }
}
\end{lstlisting}
\section*{Total Token Usage}
\textbf{Input tokens: }322
\\\textbf{Output tokens: }1023
\\\textbf{Reasoning tokens: }871
\\\textbf{Sum of `total tokens': }1345
\section*{Experiment Timings}
\textbf{Overall Experiment} started at 1780655869708, ended at 1780655904330, lasting 34622ms (34.62 seconds)
\\\textbf{Iteration \#1} started at 1780655869710, ended at 1780655904330, lasting 34620ms (34.62 seconds)
\end{document}
